\section{Introduction to Descriptive Statistics}

\emph{Statistics} is the science of collecting, organizing, analyzing, and interpreting data to make informed decisions. It is divided into two main branches:

\begin{itemize}
	\item \textbf{Descriptive statistics:} Focuses on summarizing and describing the characteristics of a dataset using numerical measures and graphical representations.
	\item \textbf{Inferential statistics:} Uses a sample of data to make generalizations and predictions about a larger population.
\end{itemize}

In this chapter, we will focus on descriptive statistics, which serves as the foundation for understanding more advanced concepts in probability and statistical inference.

\subsection{Why is descriptive statistics important?}

In the era of data, the ability to summarize and visualize information is fundamental. Descriptive statistics allows us to:

\begin{itemize}
	\item \textbf{Summarize large volumes of data} into a few representative measures.
	\item \textbf{Identify patterns and trends} in the data.
	\item \textbf{Compare different datasets} objectively.
	\item \textbf{Detect outliers} that may indicate errors or interesting phenomena.
	\item \textbf{Communicate results} clearly and concisely.
\end{itemize}

\subsection{Fundamental concepts}

Before studying descriptive measures, it is important to distinguish between:

\begin{definition}[Population and sample]
	\begin{itemize}
		\item \textbf{Population:} The complete set of elements or individuals sharing a common characteristic about which we want to obtain information.
		\item \textbf{Sample:} A representative subset of the population selected for the study.
	\end{itemize}
\end{definition}

\begin{definition}[Parameter and statistic]
	\begin{itemize}
		\item \textbf{Parameter:} A numerical measure describing a characteristic of the \emph{population} (for example, the population mean $\mu$).
		\item \textbf{Statistic:} A numerical measure calculated from a \emph{sample} (for example, the sample mean $\bar{X}$).
	\end{itemize}
\end{definition}

\begin{remark}
	In practice, we can rarely study the entire population, so we use sample statistics to estimate population parameters.
\end{remark}

\subsection{Types of descriptive measures}

Descriptive measures are classified into:

\begin{enumerate}
	\item \textbf{Measures of central tendency:} Indicate the value around which data tend to cluster (mean, median, mode).
	\item \textbf{Measures of dispersion:} Indicate how spread out the data are (range, standard deviation, variance).
	\item \textbf{Measures of position:} Indicate the relative position of a value within the dataset (quartiles, percentiles).
	\item \textbf{Measures of shape:} Describe the shape of the data distribution (skewness, kurtosis).
\end{enumerate}

In the following sections, we will study the measures of central tendency and dispersion, which are the most widely used in practice.
